\section{Introduction}

\subsection{Enabling Safe BVLOS and UAM Operations}

The rapid expansion of Beyond Visual Line of Sight (BVLOS) drone operations and Urban Air Mobility (UAM) promises transformative economic and societal benefits, ranging from autonomous parcel delivery to passenger air taxis \cite{faa2023uam}. However, safe integration into dense urban airspace requires overcoming critical infrastructure challenges. Central among these is continuous, reliable surveillance coverage to support Detect-and-Avoid (DAA) capabilities for both manned and unmanned aircraft \cite{mueller2017enabling}.

Current manned air traffic relies on cooperative transponders and centralized Air Traffic Control (ATC), but BVLOS/UAM operations introduce uncooperative airspace users—malfunctioning drones, hobbyist UAVs, or unauthorized vehicles—that must still be detected to prevent collisions and ensure public safety \cite{faa2020utm}. The resulting surveillance uncertainty directly limits airspace capacity: regulators may mandate conservative separation minima (e.g., 500 m radius) when detection reliability is low, whereas high-confidence monitoring can reduce separations by an order of magnitude \cite{nasa2021utm}. This makes surveillance infrastructure a critical enabler of operational scalability for urban air corridors.

\subsection{Ground-Based Surveillance Systems (GBSS)}

Ground-Based Surveillance Systems (GBSS) provide the foundational monitoring layer for UTM, complementing onboard DAA sensors with persistent infrastructure-supported detection \cite{kopardekar2016unmanned}. Multiple sensor modalities are employed, each with trade-offs:

\begin{itemize}
    \item \textbf{Radar}: robust in adverse weather and all lighting conditions, but limited in detecting small radar cross-section (RCS) UAVs at longer distances \cite{guvenc2018detection,semkin2020analyzing}.
    
    \item \textbf{RF/RID Receivers}: cost-effective and low-power, but dependent on cooperative transmitters \cite{guvenc2018detection,besada2022cuas}.
    
    \item \textbf{Electro-Optical/Infrared (EO/IR) sensors}: high-resolution visual coverage, but degraded by fog, rain, or night conditions \cite{guvenc2018detection,yan2023passive}.
    
    \item \textbf{Acoustic sensors}: can detect propeller noise without line-of-sight, but are highly susceptible to urban background noise \cite{guvenc2018detection,yan2023passive}.
\end{itemize}

This diversity motivates heterogeneous multi-sensor networks, combining complementary modalities to ensure robust coverage and redundancy, as emphasized in counter-UAS and UTM surveillance literature \cite{muscat2023,besada2022cuas,blasch2021information}. Despite significant interest, designing cost-effective GBSS remains challenging: one must optimize the number, type, and placement of sensors over large urban areas, balancing coverage, redundancy, and budget while considering deployment constraints such as rooftop accessibility, line-of-sight obstructions, and regulatory limits \cite{aievola2022ground,dulia2024sand}.

\subsection{Prior Work on Sensor Placement Optimization}

Several frameworks have addressed GBSS or C-UAS sensor allocation (detailed comparison in Section~\ref{sec:related}, Table~\ref{tab:related_comparison}):

\textbf{MUSCAT} \cite{muscat2023}: multi-sensor C-UAS placement with coverage/gap/cost metrics and exhaustive search over elevation-oriented terrain models. Rigorous for small candidate sets, but $O(2^n)$ search is intractable beyond roughly 10--15 locations.

\textbf{Radar / AAM network design}: Aievola et al.\ analyze ground-based radar networks for UAM \cite{aievola2022ground}; Dulia and Shihab optimize AAM surveillance networks under cost and detection-probability floors over terrain classes \cite{dulia2024sand}. These advance placement for cooperative/non-cooperative AAM traffic but do not jointly expose dual-layer OSM+ray-trace CapEx Pareto fronts of the form pursued here.

\textbf{Metaheuristics}: NSGA-II and related MOEAs are standard for coverage--cost placement in WSN and radar management \cite{deb2002fast,molina2008layout,iqbal2015wsn,charlish2016radar}, motivating transfer to occlusion-aware GBSS chromosomes.

\subsection{Remaining Challenges}

Despite progress, three gaps persist in urban GBSS design (Table~\ref{tab:related_comparison}):

\textbf{Scalability}: Exhaustive and many MIP formulations become impractical for tens of heterogeneous sensors across square kilometers with occlusion physics \cite{muscat2023,katoh1998resource}.

\textbf{Urban Environment Fidelity}: Statistical LoS models omit building-resolved shadowing that dominates low airways \cite{3gpp38901,itur_p1411}; OSM-enabled deterministic LoS for C-UAS CapEx planning remains comparatively sparse relative to telecom uses of OSM \cite{osm_applications}.

\textbf{Operational Metric Split}: Cooperative RID/RF coverage and non-cooperative dark-target coverage are often collapsed into one score, understating dark-target risk for BVLOS safety cases \cite{besada2022cuas,blasch2021information}.

\subsection{Contributions}

This paper revises and extends prior city-scale GBSS optimization (MUSCAT / SCOPAS lineage) with a dual-layer operational framing and two claim-closing controls:

\begin{enumerate}
    \item \textbf{Dual-layer multi-objective placement}: NSGA-II optimizes threat-weighted \(M_{wp,coop}\) and \(M_{wp,noncoop}\) jointly with CapEx, instead of a single fused coverage score.
    \item \textbf{Heterogeneous non-cooperative mix}: Radar, EO, and Acoustic form the dark-target layer; RF is cooperative-only. Acoustic provides a documented low-CapEx non-coop filler---without claiming universal heterogeneity dominance on classic \(M_c\).
    \item \textbf{Requirement floors for planning}: Explicit joint corridors (e.g., \(0.90/0.35\)) with Pareto filtering replace opaque GREEN thresholds and unsupported ``\% of optimum'' claims.
    \item \textbf{Evidence across scenes}: Synthetic city asset-weighted fronts, airport split coop/noncoop runs, and short target-seeking demos quantify the coop-cheap / noncoop-hard asymmetry for BVLOS safety cases.
    \item \textbf{Claim-closing controls}: (A)~small-$n$ oracle comparison (Exhaustive / Greedy / Random / GA) under the same ray-tracer; (B)~matched-budget classic \(M_c\) modality ablation; (C)~dual-layer CapEx ablation on \(M_{wp,coop}/M_{wp,noncoop}\). Together these bound GA gaps where enumeration is feasible and show when heterogeneity helps (joint dark-target floors) versus when RF-only wins (lean classic/coop coverage).
\end{enumerate}

Locked contribution sentence: NSGA-II with deterministic 3D ray-tracing and OSM geometry yields dual-layer GBSS Pareto fronts spanning cooperative \(M_{wp}\) near \(0.07\)--\(0.96\) and non-cooperative \(M_{wp}\) near \(0.00\)--\(0.56\) on a synthetic city (with complementary airport and Paulista lite fronts), in a \(g\cdot p\) evaluation regime where exhaustive MUSCAT search is intractable for \(n\gtrsim25\) sites; city-scale optimality gaps and universal heterogeneity dominance remain unsupported.

The remainder of this paper is organized as follows: Section II reviews related work and the systems comparison table. Section III formulates the dual-layer problem. Section IV--V describe methodology and SCOPAS implementation. Section VI presents experimental setup. Section VII analyzes dual-layer results and Controls~A--C. Section VIII discusses implications and limits. Section IX concludes.

